\documentclass{article}
\usepackage{pathfinder-readable}

\title{Opening Offers and Bargaining Power in Two LLM Markets}

\begin{document}
\maketitle

\begin{abstract}
Paper Q studies language-model traders in a double auction, while paper P studies their negotiated contracts
in a resource game. This thread asked whether P's apparent advantage for one player reflects a better option
outside the bargain, the chance to propose first, or both. Q reports that later quotes respond to an opening
offer, and P reports favourable point contracts for the player who always proposes first, but the papers do
not separate the two causes in P. The thread paused because that separation needs a small, repeated experiment
in which player role and proposal order vary independently.
\end{abstract}

\section{Introduction}

Q places language-model buyers and sellers in a continuous double auction. A buyer values a unit, a seller
bears a cost, and a trade occurs as soon as a bid and an ask cross. Q reports slow or absent convergence
towards market equilibrium and fewer trades than the market could support. It also reports that quotes from
agents on the other side track the first quote, even though the trading rule does not require those agents to
follow it \cite{Q}.

P places two language-model players, P-Red and P-Blue, on a board where they need coloured chips to reach a
goal. They can trade chips and negotiate contracts that enforce later transfers of chips or finishing points.
Contracts can improve cooperation when a player must first help another player and wait for a return. On P's
asymmetric boards, P-Red can reach the goal alone, while P-Blue needs help. P-Red also opens every contract
negotiation \cite{P}.

The connection pursued was a question about bargaining power. An independent player can refuse a poor deal and
still finish; a first proposer can also set the starting terms. Q's opening-quote pattern suggested that
proposal order might matter in P even when neither player has a stronger option outside the contract. That
suggestion is a reason to test, not evidence that Q's market and P's game share a causal mechanism.

\section{What was done}

The peers first asked whether P's enforceable contracts could repair Q's missed trades. They checked Q's
transaction rule against P's reason for using contracts. In Q, a crossing quote executes at once. In P, a
promise to help later may be broken unless a contract enforces it. The peers therefore set aside a direct
transfer of P's commitment remedy to Q. Q instead connects low trade volume to small successive quote changes
and agents' reluctance to cross before the round ends \cite{Q,P}.

One peer then considered whether contract-like wording could change behaviour even when Q's offers already
bind. Another proposed comparing total gains with each trader's gain over an outside option. Checks against
Q's prompts and P's measure of mutual benefit narrowed both ideas. The thread finally focused on the cleaner
question: in P's point contracts, how much of P-Red's advantage comes from being able to finish independently,
and how much from proposing first? The peers checked Q's opening-quote result, P's fixed assignment of roles
and proposer, and P's own discussion of first-mover advantage \cite{Q,P}.

The peers also made limited prior-work searches. The ledger records studies of anchoring in language-model
negotiation \cite{Takenami2025,Rios2025} and adjacent work on binding offers, natural-language contracts, and
supply-chain bargaining \cite{BindingOffers2026,LanguageContracts2026,SupplyChain2026}. Those searches did not
establish that the proposed experiment is novel.

\section{Results}

The transaction rules give a narrow result about the first idea. Let \(v\) be a buyer's value for one unit,
\(c\) the seller's cost, and \(p\) the price of an immediate trade. Their payoffs are \((v-p,\,p-c)\). If a
contract merely transfers \(t\) from the buyer to the seller when that same trade occurs, the payoffs become
\[
  (v-p-t,\,p+t-c).
\]
These are exactly the payoffs from an immediate trade at the effective price \(p+t\). Under the peers' stated
assumptions of one unit, feasible prices, and no later or third-party obligations, the transfer adds no new
payoff possibility. P's contracts solve a different problem: they make later performance credible. This
argument does not rule out a change in trading behaviour caused by the wording of an offer \cite{Q,P}.

For the main question, P reports that P-Red receives about \(60\%\) of negotiated points in its programmatic
points contracts even on independent and mutually dependent boards, where P-Red lacks the asymmetric
outside-option advantage. P calls this a possible first-mover advantage. On asymmetric boards, however, P-Red
is both the independent player and the first proposer. Its favourable contract terms there cannot be divided
into a role effect and an order effect from the reported comparisons. P's finding that full state visibility
improves P-Red's terms on asymmetric boards leaves the same distinction unresolved \cite{P}. Q's response to
opening quotes makes an order effect plausible, but it supplies no estimate of that effect in P \cite{Q}.

\section{What did not hold}

The peers considered an objection based on individual losses hidden by aggregate gains. Q's exchange permits
loss-making offers, but Q's prompts tell buyers not to bid above their value and sellers not to ask below
their cost. If agents obey those instructions, an executed trade gives each side a non-negative gain. Q does
not report how often agents disobey them, so the supplied material does not show a widespread loss problem
\cite{Q}. P's measure requiring both players to beat their baseline is strict: it also fails when one player
merely matches that baseline. A zero on that measure alone does not prove that either player was harmed
\cite{P}.

The other proposed link also remained open. P reports less subsequent trading with natural-language tile
contracts than with programmatic tile contracts, despite similar numbers of covered moves. Q reports
opening-quote responses in a market with immediate settlement \cite{P,Q}. These are observations in different
games. They do not show that contract language would reduce trades in Q, or that cognitive anchoring caused
P's point-contract split. Likewise, an early proposal and a favourable final deal can move together without
proving why the agents chose those terms.

\section{Why the thread stopped}

The verifier paused the thread because the main connection is still a hypothesis. Q's opening-quote pattern
does not establish an effect of proposal order on P's contract bargains, and P already notes the possible
first-mover advantage. The supplied work contains no outcomes from a test that separates order from the
independent player's outside option.

The smallest restart is to run P's programmatic points contract on matched original and red/blue-swapped
asymmetric boards, while independently randomising which player proposes first and repeating matched runs.
Keep the model pairing, board geometry, negotiation limit, and judge fixed. Compare net point transfers and
final rewards to the independent player across the four role-by-order combinations; also record agreement
rates, joint reward, gains over baseline, and opening proposals. An independent-player premium when that
player proposes second would support an outside-option effect. A proposer premium within either role would
support an order effect. Opening proposals would help describe the path to agreement, but would not by
themselves prove cognitive anchoring.

\bibliographystyle{plainurl}
\bibliography{references}

\end{document}
